\scanpage{14}
\uline{Prop} Soit $X$ un schéma sur corps $k$, $k_1$ une ext. alg. close de $k$. Cond. équiv.

a) $\unsure{\forall|\deleted{V}\forall}$ $k$-schéma intègre $Y$, $X\times_k Y$ intègre.

b) $\forall$ \unsure{ext.|lxt.} $k'$ de $k$, $X_{k'}$ intègre

c) \deleted{$\unsure{\forall|\forall k}$} $X_{k_1}$ intègre

\uline{Déf} On dit alors que $X$ est géom intègre ($\Longleftrightarrow$ géom irréd et géom réduit)

Ainsi géom intègre $\Rightarrow$ intègre, réciproque vraie ssi $k$ alg clos.

$X_k/k$ géom int. $\Longleftrightarrow X_{k'}/k'$ géom. intègre

Stabilité par produits

\par\noindent\rule{.22\linewidth}{.4pt}\par
\uline{Prop} Soit $X$ sur corps $k$, $Z\xrightarrow{u}X$, $X$ connexe \inserted{\deleted{(irréductible)}}, $Z$ géom. connexe \inserted{non vide \deleted{(resp. géom irréd)}}, (p. ex \inserted{non vide} radiciel$/k$).

\noindent\vrule\enspace Alors $X$ est \deleted{\unsure{géom.|géomt}} connexe \deleted{(géom irréd)}.

Soit $k'$ \deleted{la clôture sép. de $X$.} \inserted{extension finie galoisienne}. Alors $\Pi=\operatorname{Gal}(k'/k)$ opère sur $X_{k'}=X\times_k\operatorname{Spec}(\unsure{k'|k'_i})$ et sur $Z_{k'}$, de façon à commuter à $u_{k'}: Z_{k'} - X_{k'}$.

\deleted{Les} \inserted{Soit $U$} parties ouvertes et fermées de $X_{k'}$, \deleted{\unsure{si il est|s'il est}} \inserted{prouvons qu'elle} \deleted{\unsure{qu'on, considérons|est $\varnothing$, considérons}} \inserted{est $\varnothing$ ou $X_k$} prouvons qu'elle est stable sous l'action de $\Pi$. En effet,

\deleted{\unsure{Si|Soit} $g\in$ $U$ contient} Quitte à rempl. $U$ par $\complement U$, on peut supposer que $U$ contient $u_{k'}(Z_{k'})$, qui est stable par $\Pi$, donc l'intersection \inserted{\unsure{$\forall$|de|$\pi$}} des

\deleted{ouverts fermés qui contiennent $u_{k_1}(Z_{k'})$ est \deleted{\unsure{bien|un}} stable sous $\Pi$,} \unsure{transformés,|transformés} de $U$ par les $g\in \unsure{\Pi|\mathrm{G}}$ est un ouvert stable \deleted{tel que} : \deleted{donc c'est l'image inverse d'une partie de $X$} non vide, provenant donc d'un ouvert \deleted{\unsure{et fermé|stable de}} de $X$, etc

\par\noindent\rule{.94\linewidth}{.4pt}\par
\uline{Ex} Composante neutre \inserted{$G^{\unsure{0|\circ}}$} d'un groupe algébrique \deleted{[elle est quasi-compacte} [Quand elle est ouverte, p. ex. $G$ loc. de t. f. sur $S$, alors on la munit de la structure induite]

$G_{\mathrm{red}}$ pour $k$ \uline{parfait} : un sous-groupe fermé de $G$

$G^{0}_{\mathrm{red}}$ id

[Soit $U$ \deleted{\unsure{et|deux}} \inserted{$\unsure{\varnothing|V}$} ouverts de $G$ ; alors $U.\unsure{U|V}^{-1}\supset G^0$ [\,$\Rightarrow$ et si $U,V$ \unsure{contient|contiennent} l'\unsure{élémt|élément} neutre, $UV\supset G^0$] [Cor. $G^0$ quasi-compact.]

% archive sheet number: 13
